\emph{Centera uses cryptographic hash functions as a means of addressing
stored objects, thus creating a new class of data storage referred to as
CAS (content addressed storage). Such hashing serves the useful function
of providing a means of uniquely identifying data and providing a global
handle to that data, referred to as the Content Address or CA. However,
such a model begs the question: how certain can one be that a given CA
is indeed unique?}

\emph{In this paper we describe fundamental concepts of cryptographic
hash functions, such as collision resistance, preimage resistance, and
second-preimage resistance. We then map these properties to the MD5 and
SHA-256 hash algorithms, which are used to generate the Centera content
address. Finally, we present conditional estimates of accidental collision probability for the
Centera Content Address and distinguish them from adversarial security.}
